
\documentclass[UTF8,12pt,a4paper]{ctexart}

% 页面与排版
\usepackage{geometry}
\usepackage{setspace}
\usepackage{graphicx}
\usepackage{float}
\usepackage{booktabs}
\usepackage{tabularx}
\usepackage{enumitem}
\usepackage{xcolor}
\usepackage{hyperref}
\usepackage{xurl}
\usepackage{fancyhdr}
\usepackage{titlesec}
\usepackage{textcomp}

\geometry{
    top=2.5cm,
    bottom=2.5cm,
    left=2.8cm,
    right=2.8cm
}

\setstretch{1.25}
\setlength{\parindent}{2em}
\setlength{\parskip}{0.4em}

% 颜色
\definecolor{mainblue}{RGB}{20,78,121}

% 超链接
\hypersetup{
    colorlinks=true,
    linkcolor=mainblue,
    urlcolor=mainblue
}

\urlstyle{same}

% 标题格式
\titleformat{\section}
{\Large\bfseries\color{mainblue}}
{\thesection}
{0.8em}
{}

\titleformat{\subsection}
{\large\bfseries}
{\thesubsection}
{0.8em}
{}

% 页眉页脚
\pagestyle{fancy}
\fancyhf{}
\fancyhead[L]{南安普顿留学攻略}
\fancyhead[R]{\leftmark}
\fancyfoot[C]{\thepage}
\renewcommand{\headrulewidth}{0.4pt}

% 文档信息
\title{
    \vspace{2cm}
    {\Huge\bfseries 南安普顿留学攻略}\\[0.5cm]
    {\Large Southampton 留学生活指南}
}

\author{作者：Zhehao Qian}
\date{\today}

\begin{document}

% 封面
\begin{titlepage}
    \centering

    \vspace*{4cm}

    {\Huge\bfseries\color{mainblue} 南安普顿留学攻略\par}

    \vspace{0.8cm}

    {\Large Southampton 留学生活指南\par}

    \vfill

    {\large 作者：Zhehao Qian\par}

    \vspace{0.3cm}

    {\large 更新时间：\today\par}

    \vspace{2cm}
\end{titlepage}

% 目录
\tableofcontents
\clearpage

% 正文
\section{前言}

笔者闲来无事，加上导师迟迟不回消息，于是决定来写一篇这个。

首先，非常感谢我的留学中介老师吴悦老师。希望大家是从她手中拿到这份攻略的。

本 PDF 使用 \LaTeX{} 编写，并通过 Overleaf 的 XeLaTeX 编译。源文件之后会在 GitHub 上开源，如果有人需要修改或者二创，请移步：

\begin{center}
    \texttt{[GitHub 链接待补充]}
\end{center}

文章大纲由 ChatGPT 生成，我就把它当作一份 Q\&A，然后逐项填写内容。正文均为手写。至于一些我认为不太重要的内容，请大家自行在小红书或者 Google 上搜索。本质上，我就是一个把各种小红书帖子聚集在一起的人。

文中的很多内容都只是我自己的选择和经历，大家完全可以自行研究，并比较不同的选项。

OK，首先进行叠甲：本人 25 Fall（我 Chovy，南安你排名给我保好了啊，111 是什么雷霆排名啊），入学时南安 QS 排名第 87，目前是数据科学专业研究生在读。

\textit{注：本文出现的所有价格记录于 2026 年 6 月，仅供参考，
实际价格请以购买时的官方网站或应用为准。}

\section{认识南安普顿}

\subsection{城市概况}

南安普顿整体比较安静，居民也比较友善，安全系数很高。不过，城市里的娱乐设施相对较少，主要集中在市中心的 Westquay 附近。

南安普顿距离伦敦较近，前往其他城市也比较方便。

\subsection{南安普顿大学简介}

不想介绍了，我直接引用南安 Instagram 简介：

\begin{center}
    \textit{Among the top 100 universities in the world.}
\end{center}

不过，这里还是给出一点就读建议。

由于近期的一些排名变化，首先，理工科学生还是比较建议来读的，尤其是 ECS 学院。其次，以前均分比较低，例如 75 分左右，可能够不上世界前 100 大学的学生，现在或许也可以尝试一下。具体情况请咨询自己的老师，说不定读着读着，学校就又回到前 100 了。

\section{行前准备}

\subsection{签证与重要文件}

印象中比较重要的主要是肺结核检查和存款证明。

肺结核检查可以去距离自己最近的指定医院做。需要注意的是，尽早预约会方便很多。

存款证明方面，笔者当时存了 50 万元人民币，存满 28 天，使用的是中国农业银行。大家可以根据自己的情况，就近前往银行办理。去柜台后直接说明需要办理留学签证使用的存款证明，工作人员一般都明白。

个人建议尽量选择四大行办理。

不过，无论选择哪一种方式，都请尽可能早一点把这些事情准备好。

\subsection{行李清单}

完整的行李清单可以向自己的中介老师询问。这里主要讲一些我个人认为比较重要的事项：

\begin{enumerate}
    \item 首先是衣物。南安普顿并不是特别冷，最冷的时候穿内搭、卫衣和冲锋衣，基本上就足够了。

    如果有家长看到这段话，请记住：在这里，\textbf{防风大于保暖}。冲锋衣两到三件是刚需。

    \item 药物方面，我个人认为，只需要携带一些自己平时会吃、但又不太常见的药物。甚至心大一点的话，也可以不带太多药。

    因为这里的 Boots 和各类超市都可以买到感冒药、止泻药等常用药品，种类比较齐全，效果也不错。

    \item 重要！尽可能多带一点餐巾纸，哪怕塞在行李箱里当减震材料也好。这里的餐巾纸实在是太烂了。

    \item 可以的话，记得带一个小风扇。别问为什么，带！真的有用！

    \item 关于水质，笔者带的是复旦申花的过滤花洒，可以购买六个滤芯，然后来到这里后再收一些二手滤芯。

    复旦申花还有安装在卫生间洗手台水龙头上的过滤器，这个也比较建议购买。南安普顿的水质还是挺差的。
\end{enumerate}

\subsection{机票与入境准备}

笔者提前很久就购买了机票。

这里有一个注意事项：请尽可能选择国内航空公司的航班。国内航司如果能够经过俄罗斯领空，飞行时间通常会短很多。

同时，可以尽量选择国内转机的航班。例如，笔者是从杭州出发，选择的路线就是杭州飞北京，再从北京飞伦敦希思罗机场。

笔者也坐过一次上海直飞伦敦的航班，全程大约 15 小时，真的太坐牢了！

现在入境时似乎已经不需要填写入境卡。过海关时，直接说明：

\begin{center}
    \textit{I'm here to study at the University of Southampton.}
\end{center}

估计就可以顺利通过了。

\subsection{常用软件}

比较常用的软件包括：

\begin{itemize}
    \item Unilink：乘坐公交车；
    \item Bluestar：也可以用于查询和乘坐公交车；
    \item 各类超市的官方应用（Waitrose，sainsburys，Tesco）；
    \item TrainPal：购买火车票；
    \item National Express Coach：购买长途巴士票；
    \item Google Maps；
    \item Google Translate。
\end{itemize}

这些软件不用一次性全部研究明白，生活一段时间后，自然就会逐渐熟悉。

这里比较建议选择苹果手机，在很多场景下都会方便一点。如果使用安卓手机，也可以买一台二手 iPhone 作为备用的 iOS 设备。

\section{抵达南安普顿}

\subsection{从机场前往南安普顿}

这里需要分情况讨论。

如果你落地的是伦敦希思罗机场，那么 National Express 基本上是最好的选择。提前在上文提到的 National Express Coach 应用中购买车票，车次还算比较多。

从希思罗机场到南安普顿通常不到两个小时，票价一般不超过 \pounds{}20，具体价格取决于购票时间。

如果你落地的是 Heathrow Terminal 2 或 Terminal 3，巴士站就在附近。全程盯着 ``Central Bus Station'' 的指示牌寻找即可。

如果落地的是其他航站楼，也不用慌。机场内部有免费的交通方式。大致逻辑就是先寻找 Underground 的指示，前往 Terminal 2，然后继续寻找 ``Central Bus Station''。

当然，如果需要更加详细的路线指引，可以在小红书搜索“希思罗 National Express”，上面有非常多的图文攻略。

如果你落地的是盖特威克机场，也可以选择乘坐火车前往南安普顿。不过笔者没有实际坐过，所以这一部分请大家自行搜索小红书攻略。

如果你选择乘坐National Express来南安，那么注意，南安一共有两个公交站点，一个在市中心，另一个在学校旁边。

这里插播一条广告啊：启德能不能给点广告费啊！

笔者当年购买了启德的安护包，好像是 1.4 万还是 1.7 万元人民币，具体价格已经忘记了。服务内容包括接机、介绍学校和城市、协助办理入住，以及陪同采购等。

如果父母认为孩子第一次独自在国外生活，需要有人照应，同时也有余力支付这笔费用，那么可以选择花钱买一个安心。

\subsection{入住流程}

入住流程其实很简单：

\begin{center}
    ``Check-in, please.'' \qquad ``Passport.''
\end{center}

基本上就可以了。

不过，有些公寓的前台并不是 24 小时有人值班，所以一定要提前询问清楚。

例如，笔者居住的 Hampton Square 前台是 24 小时有人值班的，即使晚上抵达也可以办理入住。但并不是每个公寓都这样，所以记得提前找到公寓前台的邮箱并进行沟通。

如果你需要提前购买被芯、枕芯之类的物品，并将其快递到公寓，也需要提前确认前台是否接收快递。

对的，还真的有公寓前台不接收快递。

\section{住宿}

这里只简略介绍学校宿舍和校外学生公寓的一些选择逻辑。

\subsection{学校宿舍}

详情请查看南安普顿大学住宿官网：

\begin{center}
    \url{https://www.southampton.ac.uk/student-life/accommodation}
\end{center}

\subsection{校外学生公寓}

笔者居住的是 Hampton Square 的 Platinum Ensuite（现在好像没有了），价格为每周 \pounds{}237，非常不值。

建议大家如果选择 Ensuite，可以考虑每周 \pounds{}200 左右的房间；或者直接把预算提高到每周 \pounds{}270 左右，选择 Studio。

选择公寓时，主要需要考虑楼层和采光。个人建议选择三四层及以上、朝南的房间。

南安普顿阳光比较充足，天气也很好，所以朝南是一件非常舒服的事情。

至于安静问题，大部分公寓其实都挺安静的，不需要过于担心。

\subsection{租房注意事项}

如果选择 Ensuite，就要使用公共厨房。

说实话，你大概率碰不到特别正常的人。笔者就碰到了一个超级无敌脏的本科生，属于感觉完全没有生活自理能力的那种人。

所以，如果你准备住 Ensuite，就需要提前做好公共区域可能比较脏的心理准备。

感谢GPT，大家还需要注意
\begin{enumerate}
    \item 合同是否包含水电网。
    \item 入住当天要拍照。
    \item 提前退租是否可行,合同结束前如何续租。
    \item 房间损坏如何报修。
\end{enumerate}

\section{交通}

\subsection{公交}

使用学校邮箱注册 Unilink 账户后，直接在应用中购买车票即可。

票种包括次卡、日卡、团体票、月卡、季卡和年卡。

其中，团体票是我自己的叫法。如果父母来看你，可以购买这种票，一天大约 \pounds{}9，三个人可以无限次乘坐公交车。

确定一整年都需要乘坐公交车上下学的同学，可以直接购买年卡，价格大约为 \pounds{}345。

平时乘坐公交车比较少的同学，可以购买次卡，单次价格大约为 \pounds{}1.50。

\subsection{火车}

南安普顿常用的火车站主要有两个：Swaythling Station 和位于市中心的 Southampton Central Station。

住在学校附近的同学，通常使用 Swaythling Station 会比较方便。

火车票可以在 TrainPal 上购买。记得提前办理 Railcard，可以节省不少火车票费用。

\section{日常生活}

\subsection{超市与购物}

南安普顿市中心有 ASDA、M\&S 和 Waitrose，其中市中心的 Waitrose 规模比较小。

Portswood 有一家规模较大的 Sainsbury's 和一家中等规模的 Waitrose。

因为笔者没有去过 Lidl 和 Aldi，所以这里就不进行介绍。

商品品质方面，个人感觉 Waitrose 大于 Sainsbury's。不过，这两家基本上已经能够满足大部分日常购物需求。

中超方面，我基本上只去 Portswood 的小恐龙，它家的奶茶也可以。不过，南安普顿的中超还是挺多的，大家可以自己到处逛逛。

这里不进行任何具体商品推荐。

小红书上有非常多商品推荐，但说实话，很多东西在小红书上被吹得天花乱坠，我自己吃起来却觉得非常一般。这种情况比比皆是，所以还是需要自己尝试以后再判断。

\subsection{餐饮}

这里稍微做一些餐厅推荐，不过本人在外面吃过的餐厅并不是特别多。

中餐方面，市中心可以选择上海滩和甜森，个人评价是甜森 $>$ 上海滩。

学校附近的 Burgess Road 上，食记、风味面馆和港式茶餐厅都挺不错。此外，还有距离学校比较近的车仔档和食满家。个人觉得一品居挺一般的，可能是个人口味问题。

日料方面，可以选择市中心的千都，以及 Portswood 的 UMAMI。

这里重点推荐一下 Hampton Square 楼下的 Nathons。价格不贵，汉堡也很好吃，我已经吃过无数顿了，无敌。

\subsection{银行卡}

本人选择的是工商银行星座卡和汇丰银行卡。

南安普顿市中心有一家汇丰银行的线下网点。具体办理流程需要等你抵达南安普顿，拿到公寓地址和邮箱钥匙后，再前往汇丰银行官网申请。

现在汇丰办卡的流程非常简单，简单到我已经不记得自己当时提供了哪些材料。

其他选择还包括小马卡。日常买菜也可以办理一张 Monzo 卡，好像会有返现活动。此外，还可以选择一些虚拟银行卡。

其实，单纯使用国内发行的 Visa 信用卡也完全够用。只不过办理一张英国本地银行卡，在一些场景下会方便一点。

\subsection{电话卡与网络}

笔者只使用过 CMLink，整体信号还不错，下飞机一下就有信号了。

CMLink支持提前邮寄到你的家里，并且预定一个激活日期。这个我觉得还是很方便的。

我使用的是每月 \pounds{}9、包含 20GB 流量的套餐，日常基本上够用。

我记得黑色星期五期间好像有充值 \pounds{}50 赠送 \pounds{}15 的活动。笔者现在是 2026 年 6 月，去年 12 月充值的钱到现在还没有用完。

记得可以在出发前用英国的手机号注册一个新的iOS账户，然后到了英国这里之后进行转区才可以下载上面说的一些APP。

\section{医疗与安全}

\subsection{GP 注册与看病}

GP 注册在距离自己较近的地方即可。

例如，笔者住在学校附近的 Hampton Square，所以就注册在学校医务室。

呵呵，这里请注意：在来英国之前，一定要好好检查自己的牙齿，以及其他可能需要处理的疾病。

虽然你已经支付了 NHS 相关费用，也就是医保费用，但是这里支持医保的公立医疗服务，拔一颗牙可能需要预约四周。

四周？我牙都痛死了。

如果前往私人诊所，拔一颗牙可能需要 \pounds{}700，和直接抢钱差不多。

\subsection{日常安全与诈骗}

日常生活方面，南安普顿其实挺安全的。

关于诈骗，我的处理方式也很简单：不接任何陌生电话，只注意查看邮件。

笔者到目前为止无事发生。

\section{生活费用}

笔者每个月的日常生活费用大约为 \pounds{}700。不点外卖，但是大部分时间都会在外面吃饭。

住宿一共花费大约 11 万元人民币，这一部分其实还可以节省一些。

学费大约为 33 万元人民币。不同专业的学费会有差异，我读的是数据科学专业，本身就比较贵。

雅思、机票和其他杂七杂八的费用，可以暂时按照 5 万元人民币计算。笔者学习雅思时没有报班，如果需要报班，则需要额外计算费用。

一年的生活费大约为 8--9 万元人民币。

如果全部自己做饭，理论上可以将每月生活费用控制在 \pounds{}200--\pounds{}400 之间。不过，那样生活会比较累，所以还是需要根据自己的情况进行判断。

每个人的生活方式都不同，我的数据只能作为参考。

\section{一些生活小贴士}

这里主要记录一些我认为在日常生活中比较有用的东西，想到什么就说什么。

\begin{enumerate}
    \item 关于回国加速器，笔者没有选择比较常用的 QuickFox，而是选择了 VK 加速器全球版，学生可以享受优惠。

    不过，实际使用效果比较一般，所以还是更建议选择 QuickFox。

    \item 这里的空气还是有一点干，可以携带一个小型加湿器，或者抵达英国以后再购买一个。

    \item 这里的下水道比较垃圾，洗澡时使用的地漏真的很容易彻底堵死。

    使用威猛先生的下水道疏通产品会非常有效，也可以直接向公寓报修。

    \item 关于喝水，笔者平时喝的是 Volvic 矿泉水。

    其他品牌的水喝得比较少，但我总感觉会有一种硬水的味道。购买矿泉水每个月大约需要花费 \pounds{}20。可以叫Sainsburys的外卖，用Saver Slot 1p送到家。

    也可以选择 BRITA 滤水壶，据说过滤后的水非常干净。

    另外请注意：只有厨房里的冷水是直饮水。卫生间洗手池里的冷水也不建议直接饮用，热水则更是绝对不能直接饮用。
\end{enumerate}

\section{总结}

可以看出，这篇攻略整体使用的都是比较日常、比较口语化的表达方式，本质上就是抱着一种聊天的心态，为大家解答一些问题。

文中出现了大量“小红书”这个词。来到英国留学以后，小红书基本上就是你的百度百科。你会发现，上面真的什么内容都有。

我的这篇攻略更多是为了让你知道“原来还有这么一个东西”，然后方便你根据关键词继续进行搜索。

如果真的要把所有事情都事无巨细地解释清楚，这篇攻略可能会写到和毕业论文差不多长。

至于为什么没有写学校学习的部分，我实在是不知道该写什么，这里给大家推荐一个xhs博主‘南安今天可心了吗’（xkxxkxxkxzka），她的帖子什么都有。

另外，笔者本人也不太喜欢出门玩，所以大家会发现，文中几乎没有介绍娱乐和旅行相关的内容。

这一部分就留给大家自行发掘吧。毕竟玩这种事情，总是能够玩明白的。

最后，祝大家留学顺利！

% 图片示例
% \begin{figure}[H]
%     \centering
%     \includegraphics[width=0.8\textwidth]{images/example.jpg}
%     \caption{图片标题}
%     \label{fig:example}
% \end{figure}


\end{document}

